\documentclass[../Main.tex]{subfiles}

\begin{document}
\section{Definition and Examples}
\begin{definition}{Bifurcation}
    Let $\dvec{x} = \vec{f}(\vec{x};\mu)$ be a dynamical system parameterised by $\mu$. A \underline{bifurcation} is a change in the topological structure of the flow as $\mu$ passes through some critical value $\mu_c$. In this case, $\mu_c$ is an \underline{bifurcation point}.
\end{definition}
\begin{remark}
    We will often take $\mu_c = 0$, after possibly shifting it in the definition of $\vec{f}$.
\end{remark}
\begin{example}
    Consider:
    \begin{equation*}
        \dot{x} = \mu x + y - x(x^2 + y^2),\qquad \dot{y} = -x + \mu y - y(x^2 - y^2)
    \end{equation*}
    and transforming this to plane polar coordinates gives:
    \begin{equation*}
        \dot{r} = r(\mu - r^2),\qquad \dot{\theta} = 1.
    \end{equation*}
    We can consider the fixed point at the origin for different values of $\mu$: for $\mu < 0$ we have a stable focus, for $\mu = 0$ we still have a stable focus but slower convergence, and for $\mu > 0$ we have an unstable focus with periodic orbit. Therefore, we conclude that there is a bifurcation at $\mu = 0$.
\end{example}
\begin{remarks}
    \item We will generally be concerned with local bifurcations, looking at topological changes in a small vicinity of fixed points.
    \item We can say that if a fixed point changes stability then there has been a change in topological structure.
        Therefore local bifurcations happen at non-hyperbolic fixed points.
\end{remarks}
\begin{proposition}
    Local bifurcations require the real part of at least one eigenvalue of a fixed point to be zero.
    \label{propBifZeroEVal}
\end{proposition}
\begin{proof}[sketch]
    We can linearise near the fixed point and apply \thmref{propHypMeansLinearOnly} to show that the topological structure changes only when the fixed point becomes non-hyperbolic at the bifurcation point.
\end{proof}
\section{Bifurcations in One Dimension}
\subsection{Saddle-Node and Transcritical Bifurcations}
Consider a 1D system, $\dot{x} = f(x; \mu)$, with a non-hyperbolic fixed point at $x = 0$ when $\mu = 0$. This means $f(0, 0) = 0$ and $\frac{df}{dx}(0, 0) = 0$.

This is related to stationary bifurcations in $\R^n$, where at least one eigenvalue is exactly zero (not just its real part).
\begin{example}
    Consider the dynamical system $\dot{x} = \mu - x^2$. We see that for $\mu < 0$, $\dot{x}$ is always negative, and for $\mu = 0$ we only have the origin where $\dot{x} = 0$. When $\mu > 0$, we will have a region centred around the origin for which $\dot{x} > 0$. See figure~\ref{figOneDBifs} for these cases.
\end{example}
\begin{figure}
    \centering
    \begin{tikzpicture}[scale=1]
        \draw (-2, 2) -- (2, 2)
            node[pos=0, left] {$\mu < 0$}
            pic[pos=0.3, rotate=90] {patharrow}
            pic[pos=0.7, rotate=90] {patharrow};

        \draw (-2, 0) -- (2, 0)
            node[pos=0, left] {$\mu = 0$}
            pic[pos=0.3, rotate=90] {patharrow}
            pic[pos=0.7, rotate=90] {patharrow};
        \fill (0, 0) circle[radius=0.1];

        \draw (-2, -2) -- (2, -2)
            node[pos=0, left] {$\mu < 0$}
            pic[pos=0.2, rotate=90] {patharrow}
            pic[pos=0.5, rotate=-90] {patharrow}
            pic[pos=0.8, rotate=90] {patharrow};
        \fill (-0.5, -2) circle[radius=0.1];
        \fill (0.5, -2) circle[radius=0.1];
    \end{tikzpicture}
    \caption{Diagram of 1D bifurcations}
    \label{figOneDBifs}
\end{figure}

We can expand $f(x;\mu)$ near $x = 0, \mu = 0$:
\begin{align*}
    f(x;\mu) &= f(0;0) + \mu f_\mu(0, 0) + \frac12 f_{xx}(0;0)x^2 + x\mu f_{x\mu}(0;0) + \frac12 f_{\mu\mu}(0;0) + \cdots \\
    &= \left[\mu f_{\mu}(0; 0) + O(\mu^2)\right] + x\mu f_{x\mu}(0; 0) + \frac12 f_{xx}(0;0) + O(x^3) \\
    &= A\mu + B x\mu + Cx^2 + O(\mu^2, x^3)
\end{align*}
where $A = f_\mu(0;0), B = f_{x\mu}(0;0), C = \frac12 f_{xx}(0;0)$. Completing the square:
\begin{equation*}
    \dot{x} = A\mu + C\left(x + \frac{B\mu}{2C}\right)^2 - \frac{B^2 \mu^2}{4C} + O(\mu^2, x^3)
\end{equation*}
but we can remove this last term because we are working to $O(\mu)$:
\begin{equation*}
    \dot{x} \approx A\mu + C\left(x + \frac{B\mu}{2C}\right)^2.
\end{equation*}
Transforming coordinates as:
\begin{equation*}
    x \mapsto -y + \frac{B}{2C}\mu,\qquad t \mapsto T / C, \qquad \mu \mapsto \tilde{\mu} = -\frac{A}{C}\mu
\end{equation*}
yields the \underline{normal form}:
\begin{equation*}
    \dot{y} = \tilde\mu - y^2.
\end{equation*}
This requires $A$ and $C$ to be non-zero. This type of bifurcation is known as a \underline{saddle-node} bifurcation. See figure~\ref{figSaddleNodeBif}

\begin{figure}
    \centering
    \begin{tikzpicture}[scale=1]
        \draw[->] (-2, 0) -- (2, 0) node[right] {$\mu$};
        \draw (0, -2) -- (0, 2) node[above] {$x$};

        \draw[blue] plot[domain=0:1.4142, samples=50, smooth] ({\x*\x}, {\x});
        \draw[dashed, blue] plot[domain=0:1.4142, samples=50, smooth] ({\x*\x}, {-\x});
    \end{tikzpicture}
    \caption{Diagram for a saddle-node bifurcation}
    \label{figSaddleNodeBif}
\end{figure}

If in fact $A = 0$ we have a \underline{transcritical bifurcation}, with normal form given by:
\begin{equation*}
    \dot{y} = \tilde{\mu} y - y^2
\end{equation*}
under the transformations $y = -x, T = Ct, \tilde\mu = \frac B C \mu$. This gives a stable fixed point at the origin (and an unstable one at $x = \mu$) for $\mu < 0$, a half-stable, half-unstable fixed point at the origin for $\mu = 0$, and an unstable fixed point at the origin (and stable one at $x = \mu$) when $\mu > 0$. We can represent this as a plot of $\mu$ against $x$ -- see figure~\ref{figTranscritBif}. Here the dotted line represents an unstable fixed point, the solid line represents a stable one.
\begin{figure}
    \centering
    \begin{tikzpicture}[scale=1]
        \draw[->] (-2, 0) -- (2, 0) node[right] {$\mu$};
        \draw (0, -2) -- (0, 2) node[above] {$x$};

        \draw[dashed, blue] (-2, -2) -- (0, 0);
        \draw[blue] (-2, 0) -- (0, 0);

        \draw[dashed, blue] (0, 0) -- (2, 0);
        \draw[blue] (0, 0) -- (2, 2);
    \end{tikzpicture}
    \caption{Diagram for a transcritical bifurcation}
    \label{figTranscritBif}
\end{figure}

\subsection{Pitchfork Bifurcations}
If the dynamical system has a symmetry, and is unchanged under $x \mapsto -x$, then our previous expansion for $\dot{x}$ requires $A = C = 0$. To see the dynamics of the bifurcation we go to higher order in $x$ and $\mu$,
\begin{equation*}
    \dot{x} = \mu x f_{x\mu}(0;0) + \frac16 x^3 f_{xxx}(0;0) + \cdots
\end{equation*}
which, after a change of variables, has normal form:
\begin{equation*}
    \dot{x} = \mu x - ax^3
\end{equation*}
for $a = \pm 1$. In the case of $a = 1$, we have a \underline{supercritical pitchfork bifurcation}. This has fixed points at $x = \pm \sqrt{\mu}$ as well as $x = 0$.

In the case of $a = -1$, we have a \underline{subcritical pitchfork bifurcation}, with fixed points at $\pm \sqrt{-\mu}$. For both types of pitchfork, see figure~\ref{figPitchforkBifs}.
\begin{figure}
    \centering
    \begin{tikzpicture}[scale=1]
        % Supercritical pitchfork
        \begin{scope}
            \draw[->] (-2, 0) -- (2, 0) node[right] {$\mu$};
            \draw (0, -2) -- (0, 2) node[above] {$x$};

            \draw[blue] (-2, 0) -- (0, 0);
            \draw[dashed, blue] (0, 0) -- (2, 0);

            \draw[blue] plot[domain=0:1.4142, samples=50, smooth] ({\x*\x}, {\x});
            \draw[blue] plot[domain=0:1.4142, samples=50, smooth] ({\x*\x}, {-\x});

            \node[below] at (0, -2) {Supercritical};
        \end{scope}

        % Subcritical pitchfork
        \begin{scope}[shift={(6, 0)}]
            \draw[->] (-2, 0) -- (2, 0) node[right] {$\mu$};
            \draw (0, -2) -- (0, 2) node[above] {$x$};

            \draw[blue] (-2, 0) -- (0, 0);
            \draw[dashed, blue] (0, 0) -- (2, 0);

            \draw[dashed, blue] plot[domain=0:1.4142, samples=50, smooth] ({-\x*\x}, {\x});
            \draw[dashed, blue] plot[domain=0:1.4142, samples=50, smooth] ({-\x*\x}, {-\x});

            \node[below] at (0, -2) {Subcritical};
        \end{scope}
    \end{tikzpicture}
    \caption{Diagram of supercritical and subcritical pitchfork bifurcations}
    \label{figPitchforkBifs}
\end{figure}
\subsection{Structural Stability of One-Dimensional Bifurcations}
We will consider perturbed dynamical systems of the form:
\begin{equation*}
    \dot{x} = f(x;\mu) + \epsilon(x)
\end{equation*}
where we can assume that $f$ is in normal form.

Consider a saddle node, with $\dot{x} = \mu - x^2$.
\begin{itemize}
    \item If we add $\epsilon$ to the right-hand side,
        \begin{equation*}
            \dot{x} = \epsilon + \mu - x^2
        \end{equation*}
        which we can see does not change the dynamics (if we redefine $\mu$).
    \item If we add $\epsilon x$ to the right-hand side,
        \begin{equation*}
            \dot{x} = \mu + \epsilon x - x^2
        \end{equation*}
        and we can complete the square, rescale $x$ and redefine $\mu$ to get back to the same normal form.
\end{itemize}
Since our expansion worked to order $O(x)$, we conclude that the saddle-node is structurally stable.

Consider instead a transcritical bifurcation, with normal form $\dot{x} = \mu x - x^2$. If we perturb this for small $\epsilon$, we get $\dot{x} = \epsilon + \mu x - x^2$.
\begin{itemize}
    \item If $\epsilon < 0$, we find a region of $\mu$ around $0$ where there are no fixed points.
    \item If $\epsilon > 0$, we find that the fixed points never change stability and there is no bifurcation at all!
\end{itemize}

Now if we have a pitchfork bifurcation, with normal form $\dot{x} = \mu x - x^3$, we also find it is structurally unstable under adding a constant $\epsilon$.
\section{Bifurcations in Two Dimensions}
Consider the generalisation of what we were doing before,
\begin{equation*}
    \dvec{x} = \vec{f}(\vec{x};\mu)
\end{equation*}
We can have four possibilities depending on the form of the linearised matrix $A$ around the origin.
\begin{itemize}
    \item One eigenvalue is zero -- a \underline{simple bifurcation} or \underline{stationary bifurcation} in 2D, and it is essentially the same as our one-dimensional bifurcations.
    \item The eigenvalues are imaginary -- $\lambda = \pm i\omega$ -- an \underline{oscillatory bifurcation} or \underline{Hopf bifurcation}.
    \item Both eigenvalues are zero and non-degenerate -- a \underline{double-zero}\\\underline{bifurcation}.
    \item Both eigenvalues are zero and degenerate -- a \underline{Takens-Bogdanov} bifurcation.
\end{itemize}
The last two cases will not be covered in this course.
\subsection{The Centre Manifold}
Consider a system:
\begin{equation}
    \dot{x} = xy - \mu x,\qquad \dot{y} = -y + x^2
    \label{eqnBifSystem}
\end{equation}
There is a fixed point at the origin for any value of $\mu$, and for $\mu > 0$ we find a further two fixed points at $(\pm \sqrt{\mu}, \mu)$. Linearising around the origin gives:
\begin{equation*}
    A = 
    \begin{pmatrix}
        -\mu & 0 \\
        0 & -1
    \end{pmatrix}
\end{equation*}
and so when $\mu = 0$, the origin is a non-hyperbolic fixed point, so the dynamical system admits a bifurcation. We analyse close to the origin and bifurcation point by requiring $\abs x \sim \abs y \sim \abs \mu \ll 1$. Then $\abs{\dot x} \sim O(\mu^2)$ and $\abs{\dot y} \sim O(\mu)$, so the system evolves much faster in $y$ than $x$ towards $y = x^2$. On this curve, we find:
\begin{equation*}
    \dot{x} = x(x^2 - \mu)
\end{equation*}
and we have reduced our two-dimension problem to a 1-dimensional one. We recognise this as a subcritical pitchfork. We can formalise the above process with the following theorem. Recall the definitions of linear stable and unstable subspaces $E^S$ and $E^U$.

\begin{definition}{Linear centre subspace}
    For a fixed point, with linearisation eigenvectors $\vec{e_i}$ and linearisation eigenvalues $\lambda_i$, the \underline{linear centre subspace} is the span of eigenvectors whose eigenvalue has zero real part.
\end{definition}

\begin{theorem}[Centre Manifold Theorem]
    If the origin is a non-hyperbolic fixed point of $\dvec{x} = \vec{f}(x)$ with linear stable, unstable and centre subspaces $E^S, E^U, E^C$, then there exist stable, unstable and centre manifolds $W^S, W^U$, $W^C$ of the same dimension as the corresponding subspace, and which are tangent to the subspace at the origin and are invariant. Note that $W^S$ and $W^U$ are unique, but $W^C$ is not unique in general.
    \label{thmCentreManif}
\end{theorem}
\begin{example}[Searching for $W^C$ in $\R^2$]
    Consider:
    \begin{equation*}
        \begin{pmatrix}\dot x \\ \dot y\end{pmatrix} = \begin{pmatrix}
            0 & 0 \\
            0 & -\lambda
        \end{pmatrix}
        \begin{pmatrix}x \\ y\end{pmatrix} + \begin{pmatrix}f(x, y) \\ g(x, y)\end{pmatrix}
    \end{equation*}
    where $\lambda > 0$. Here the linearisation is done for us, and we read off $E^C = \{\vec{e_1}\}$ and $E^S = \{\vec{e_2}\}$, where here $\vec{e_i}$ are the basis vectors (as well as eigenvalues).

    We can therefore look for a manifold of the form:
    \begin{equation*}
        y = h(x) = a_0 + a_1 x + a_2 x^2 + a_3 x^3.
    \end{equation*}
    Note that $a_0 = 0$ to pass through the origin, and $a_1 = 0$ to be tangent to $E^C$. Apply invariance:
    \begin{equation*}
        \dot{y} = \frac{dh}{dx} \dot{x} 
    \end{equation*}
    and then we can substitute $\dot{y}$ and $\dot{x}$.
\end{example}
\begin{example}
    Return to the system:
    \begin{equation*}
        \dot{x} = xy - \mu x,\qquad \dot{y} = -y + x^2
        \tag{\ref{eqnBifSystem}}
    \end{equation*}
    and take $\mu = 0$. Linearising around the origin gives:
    \begin{equation*}
        A = 
        \begin{pmatrix}
            0 & 0 \\
            0 & -1
        \end{pmatrix}
    \end{equation*}
    and so we know that $E^C = \{\vec{e_1}\}$. Using $y = h(x)$,
    \begin{equation*}
        x^2 - h(x) = \frac{dh}{dx} x h(x)
    \end{equation*}
    and substituting in to fourth power:
    \begin{align*}
        x^2 - a_2 x^2 - a_3x^3 - a_4x^4 &= (2a_2 x)x(a_2 x^2) + O(x^5) \\
        x^2(1 - a_2) - a_3x^3 - a_4 x^4 &= 2a_2 x^4 + O(x^5)
    \end{align*}
    and so $a_2 = 1$, $a_3 = 0$, $a_4 = -2$.
\end{example}
\subsection{Extended Centre Manifold}
To be able to apply \thmref{thmCentreManif} to bifurcation systems, we augment $\dvec{x} = \vec{f}(\vec x)$ with $\dot \mu = 0$ and treat $\mu$ as another independent variable. This adds one more dimension to $E^C$ and $W^C$ and allows us to work near the new origin $(\vec 0, \mu = 0)$.

For example, if the centre manifold in a 2D system was $\{y = 0\}$, we would describe the extended centre manifold with the function $y = h(x, \mu)$. In this case, we would require:
\begin{itemize}
    \item $h(0, 0) = 0$ so $W^C$ passes through $(\vec 0, 0)$
    \item $h_x(0) = h_\mu(0) = 0$ so $W^C$ is tangent to $E^C$.
\end{itemize}
\begin{example}
    Return again to the dynamical system:
    \begin{equation*}
        \dot{x} = xy - \mu x,\qquad \dot{y} = -y + x^2
        \tag{\ref{eqnBifSystem}}
    \end{equation*}
    and so using the above we extend to a 3D system by adding $\dot \mu = 0$. The system becomes:
    \begin{equation*}
        \begin{pmatrix}\dot{x} \\ \dot y \\ \dot \mu\end{pmatrix} = 
        \begin{pmatrix}
            0 & 0 & 0 \\
            0 & -1 & 0 \\
            0 & 0 & 0
        \end{pmatrix}
        \begin{pmatrix}x \\ y \\ \mu\end{pmatrix} +
        \begin{pmatrix}x(y-\mu) \\ x^2 \\ 0\end{pmatrix}
    \end{equation*}
    and so let $h(x, \mu) = a_{20}x^2 + a_{11} x\mu + a_{02} \mu^2 + \cdots$. Now our equation is:
    \begin{equation*}
        \dot y = \frac{\partial h}{\partial x} \dot x + \frac{\partial h}{\partial \mu} \dot \mu = \frac{\partial h}{\partial x} \dot x
    \end{equation*}
    and we substitute in a similar process to find:
    \begin{equation*}
        h(x, \mu) = x^2 + O(x^3, \mu^3).
    \end{equation*}

    Then we are able to translate the system to one dimension, by taking $\dot{x} = -\mu x + xy$ and substituting $y = h(x)$.
\end{example}
\begin{example}
    Consider a new system:
    \begin{equation*}
        \begin{pmatrix}\dot x \\ \dot y\end{pmatrix} =
        \begin{pmatrix}\mu + x^2 + xy + y^2 \\ 2\mu - y + x^2 + xy\end{pmatrix}
    \end{equation*}
    Extending this to include $\mu$ gives:
    \begin{equation*}
        \begin{pmatrix}\dot x \\ \dot y \\ \dot \mu\end{pmatrix} =
        \begin{pmatrix}
            0 & 0 & 1 \\
            0 & -1 & 2 \\
            0 & 0 & 0
        \end{pmatrix}
        \begin{pmatrix}x \\ y \\ \mu\end{pmatrix}
         + \begin{pmatrix}x^2 + xy + y^2 \\ x^2 + xy \\ 0\end{pmatrix}
    \end{equation*}
    which means $\lambda_1 = \lambda_2 = 0$, with eigenvectors $\vec{e_1} = (1, 0, 0)^T$ and $\vec{e_2} = (0, 2, 1)^T$. This means $E^C$ is the plane $y = 2\mu$. To find $W^C$ we need $y = h(x, \mu)$ with $\frac{\partial h}{\partial x} = 0$ but $\frac{\partial h}{\partial \mu} = 2$. We therefore write:
    \begin{equation*}
        h(x) = 2\mu + a_{20} x^2 + a_{11} x\mu + a_{02} \mu^2 + O(x^3, \mu^3)
    \end{equation*}
    and substituting gives $h(x) = 2\mu + x^2 + O(x^3, \mu^3)$. The dynamics on the extended centre manifold are then given by:
    \begin{equation*}
        \dot{x} = \mu + x^2 + x(2\mu + x^2) + 4\mu^2 + O(x^3, \mu^3) = \mu + x^2 \text{ to leading order}
    \end{equation*}
    which is a saddle-node bifurcation.
\end{example}
\begin{example}[Higher dimensions -- Lorentz equations]
    Consider: %TODO: Some $b$s are actually $\sigma$s.
    \begin{equation*}
        \begin{pmatrix}\dot x \\ \dot y \\ \dot z\end{pmatrix} = \begin{pmatrix}\sigma(y-x) \\ rx - y - xz \\ -bz + xy\end{pmatrix}
    \end{equation*}
    where $b, \sigma > 0$ and $r$ are parameters. The linearisation matrix is:
    \begin{equation*}
        A =
        \begin{pmatrix}
            -b & b & 0 \\
            r & -1 & 0 \\
            0 & 0 & -b
        \end{pmatrix}
    \end{equation*}
    which corresponds to a non-hyperbolic fixed point if $r = 1$. In this case $\lambda_1 = 0, \lambda_2 = -b$ and $\lambda_3 = -(1+b)$. The centre manifold is $E^C = \{x = y, z = 0\}$.

    This suggests the change of variables $v = y-x$ to eliminate $y$, and $\mu = r - 1$ so the bifurcation is at zero. In this case, extending the dynamical system gives:
    \begin{equation*}
        \begin{pmatrix}\dot x \\ \dot \mu \\ \dot v \\ \dot z\end{pmatrix} =
        \begin{pmatrix}
            0 & 0 & b & 0 \\
            0 & 0 & 0 & 0 \\
            0 & 0 & -(1+b) & 0 \\
            0 & 0 & 0 & -b
        \end{pmatrix}
        \begin{pmatrix}x \\ \mu \\ v \\ z\end{pmatrix} + 
        \begin{pmatrix}0 \\ 0 \\ x(\mu - z) \\ x(v+x)\end{pmatrix}
    \end{equation*}
    and the two-dimensional centre subspace is $v = z = 0$. It is a 2D plane in 4D space. We look for $W^C$ which is given by:
    \begin{equation*}
        \begin{pmatrix}v \\ z\end{pmatrix} = \begin{pmatrix}V(x, \mu) \\ Z(x, \mu)\end{pmatrix}
    \end{equation*}
    It is tangent to $E^C$, which means $V$ and $Z$ have only contributions at second order and above. Note also that we only need $V$ because $\dot{x} = bv$, and so the dynamics on the space $(x, \mu)$ are given only by $V$.

    We can use perturbation theory. Assume that $V = V_2 + V_3 + \cdots$, where $V_2$ has terms of order two, and $V_3$ has terms of order 3. Similarly, assume that $Z = Z_2 + Z_3 + \cdots$. We can find $V$ and $Z$ by:
    \begin{align*}
        \dot{v} &= \frac{\partial V}{\partial x} \dot x \implies \sigma V V_x = -(1+\sigma) V + x(\mu - Z) \\
        \dot z &= \frac{\partial Z}{\partial x} \dot x \implies \sigma V Z_x = -bZ + x^2 + xV
    \end{align*}
    To second order,
    \begin{align*}
        0 &= -(1+\sigma) V_2 + x\mu \implies V_2 = \frac{x\mu}{1+\sigma} \\
        0 &= -bZ + x^2 \implies Z = \frac{x^2}{6}
    \end{align*}
    and now we do not need to go any further in $Z$, and just need $Z$ to third order:
    \begin{align*}
        \sigma V_2 \frac{\partial V_2}{\partial x} &= -(1+\sigma) V_3 - xZ_2 \\
        \sigma \frac{x\mu}{1+\sigma} \frac{\mu}{1+\sigma} &= -(1+\sigma) V_3 - x \frac{x^2}{6} \\
        \implies V_3 &= -\frac{x^3}{\sigma(1+\sigma)} - \frac{\sigma}{(1+\sigma)^3} \mu^2 x
    \end{align*}
    and we conclude:
    \begin{equation*}
        V = \frac{x\mu}{1+\sigma} - \frac{x^3}{\sigma(1+\sigma)} - \frac{\sigma}{1+\sigma} x\mu^2 + \cdots
    \end{equation*}
    and we can neglect the third term near $x = \mu = 0$. This means our evolution equation is then:
    \begin{equation*}
        \dot x = \frac{\sigma}{b(1+\sigma)} x(\mu b - x^2)
    \end{equation*}
    with normal form $X\tilde\mu - X^3$, a supercritical bifurcation.
\end{example}
\begin{remark}
    In all of the above examples, we had higher-dimensional stationary bifurcations (zero eigenvalues) which could be reduced to one-dimensional normal forms.
\end{remark}
\section{Hopf Bifurcations}
Hopf bifurcations (or oscillatory bifurcations) occur when a fixed point has eigenvalues $\lambda_{1, 2} = \pm i \omega$ where $\omega \neq 0$ when $\mu = 0$. Here we can see that $E^C$ must be at least two-dimensional.

To find the normal form of the bifurcation, consider:
\begin{equation*}
    \begin{pmatrix}\dot x \\ \dot y\end{pmatrix} = \begin{pmatrix}
        \mu & -\omega \\
        \omega & \mu
    \end{pmatrix}
    \begin{pmatrix}x \\ y\end{pmatrix} + \begin{pmatrix}f(x, y) \\ g(x, y)\end{pmatrix}
\end{equation*}
This clearly has a Hopf bifurcation at $\mu = 0$. Go to the complex plane, with $z = iy$ and $F(z, z^*) = f(x, y) + ig(x, y)$. When $\mu = 0$,
\begin{equation*}
    \dot z = i\omega z + F(z, z^*)
\end{equation*}
Expanding near $z = 0$, to second order:
\begin{equation*}
    \dot z \approx i\omega z + a_1 z^2 + a_2 zz^* + a_3 z^{*2}
\end{equation*}
Define $\xi = z + \alpha_1 x^2 + \alpha_2 zz^* + \alpha_3 z^{*2}$ such that $\dot\xi = i\omega\xi$ to second order.

Applying the same procedure to third order, we find we can eliminate all the 3rd-order terms except one term proportional to $z \abs{z}^2$. With $z$ redefined as such,
\begin{equation*}
    \dot z = i\omega z + \nu \abs{z}^2 z + \cdots
\end{equation*}
which, transforming to polar coordinates, gives:
\begin{equation*}
    \dot{r} = \Re(\nu) r^3,\qquad \dot\theta = \Im(\nu) r^2
\end{equation*}

Therefore, after doing the same manipulations as above for $\mu \neq 0$, the normal form is given (in polar coordinates) by:
\begin{equation*}
    \dot{r} = (\mu - a r^2) r,\qquad \dot\theta = \omega
\end{equation*}
where $a = \pm 1$.

In the case of $a = 1$, we find a supercritical Hopf bifurcation, and $a = -1$ corresponds to a subcritical Hopf bifurcation. For diagrams of both, see figure~\ref{figHopf}.

\begin{figure}
    \centering
    \begin{tikzpicture}[scale=1]
        \node at (0, 0) {TODO}; %TODO
    \end{tikzpicture}
    \caption{Diagram of supercritical and subcritical Hopf bifurcations}
    \label{figHopf}
\end{figure}

In a Hopf bifurcation, a single fixed point changes its stability and a periodic orbit is created ($a = 1$) or destroyed ($a = -1$) as $\mu$ crosses $0^-$ to $0^+$.
\nonexaminablesection{Bifurcations of Periodic Orbits}

\end{document}